\section{Separating instruction-header and task-block positions}
\label{app:header-distance}
The training-distance intervention moves instructions and task information
together. A separately frozen inference experiment tests whether the balanced
models retain their advantage when the generic instruction header stays
nearby and the task-bearing block moves far away. This intervenes on inference
layout; it does not separately identify the contributions of instruction and
task distance in the preceding training treatment.

\paragraph{Independent fixed slots and fresh problems.}
We draw 32 paired-parity problems with seed 2026092301 and use all six fixed
adaptation checkpoints. These public prompts have no overlap with the previous
32-problem panel, but use the previously observed held-out structure catalog.
The 24-token generic instruction header and 88-token task block move between
independent far and near slots. The latter block contains coordinate names,
equations and output-order metadata; this is not an intervention on syndrome
distance alone. The native prefix has 16,383 tokens and the answer footer
remains fixed. Unoccupied slots contain deterministic groups of whole archive
records with exactly matching token lengths. Swapping a block with its filler
preserves the complete token multiset and the other block's exact position.
All four header/task-position combinations are tested. Two additional legacy
far/near layouts retain the original serializer on the same new problems,
measuring any effect of regrouping irrelevant records in the fixed-slot design.

\paragraph{Prediction and analysis.}
Position does not change the target dependence term: $D_{\mathrm{info}}=0$.
Thus, with $h$ the header position and $t$ the task position,
\begin{equation}
 \mathrm{KL}_{\mathrm{info}}(h,t)=E(h,t),\qquad
 P(h)=E(h,\mathrm{far})-E(h,\mathrm{near}).
\end{equation}
The primary contrast is near-trained minus balanced-trained error at
$(h,t)=(\mathrm{near},\mathrm{far})$. We additionally report the task-distance
penalty $P(h)$ at each header position, header penalties at each task position,
their factorial interaction, training-arm differences and legacy discrepancies.
The predeclared stronger interpretation requires a positive primary interval,
positive direction in all three matched seeds, and balanced mean error below
$4\log2$ in both nearby-header cells for every seed. A positive arm contrast
alone need not yield a useful two-pass predictor.

All 16 target histories and one-call/information-set endpoint laws are evaluated
for every condition. Ten thousand paired problem-bootstrap draws use seed
2026092302, preserving pairing across models and layouts. Intervals remain
conditional on these six checkpoints and the selected starting lineage.
The standard numerical screen is supplemented by repeated/cross-rank checks
of all four split layouts with both all-masked and visible-history inputs.
Every cell and control is retained regardless of accuracy.

\paragraph{Results.}
All 1,152 conditions and 2,304 endpoint laws are validated. The primary error reduction is 7.6266 nats, with conditional 95\% interval [6.9081, 8.3015].
The individual adaptation-seed effects are 12.4718, 7.6350, 2.7730 nats.
The predeclared stronger-interpretation criteria are satisfied.
Table~\ref{tab:header-distance} includes every layout and both legacy controls.
\begin{table}[t]
\centering\small
\caption{Information-set KL for independent header/task positions, averaged over 32 fresh problems. FF, FN, NF and NN denote header then task position (far or near). Legacy controls use the original whole-block layouts. Seed suffixes identify the same six fixed adaptations used throughout. All entries are nats.}
\label{tab:header-distance}
\begin{tabular}{lrrrrrr}
\toprule
Model & FF & FN & NF (primary) & NN & Legacy far & Legacy near\\
\midrule
Near 11 & 13.63 & 0.000134 & 12.47 & 6.82e-05 & 9.544 & 6.704e-05\\
Balanced 11 & 1.235e-05 & 7.123e-06 & 1.344e-05 & 8.588e-06 & 4.081e-05 & 8.635e-06\\
Near 12 & 9.621 & 1.706e-06 & 7.635 & 1.017e-06 & 7.282 & 1.03e-06\\
Balanced 12 & 3.114e-06 & 3.535e-06 & 3.845e-06 & 4.145e-06 & 4.248e-06 & 4.131e-06\\
Near 13 & 2.773 & 8.744e-06 & 2.773 & 7.057e-06 & 2.773 & 7.17e-06\\
Balanced 13 & 9.141e-05 & 6.024e-05 & 0.0001524 & 0.0001064 & 8.012e-05 & 0.0001063\\
\bottomrule
\end{tabular}
\end{table}
\begin{table}[t]
\centering\small
\caption{Complete factorial contrasts within each checkpoint. $P_F,P_N$ are task-distance penalties with far/near header; $H_F,H_N$ are header-distance penalties with far/near task. $I=P_F-P_N$. All are paired mean KL differences in nats.}
\label{tab:header-factorial}
\begin{tabular}{lrrrrr}
\toprule
Model & $P_F$ & $P_N$ & $H_F$ & $H_N$ & $I$\\
\midrule
Near 11 & 13.63 & 12.47 & 1.162 & 6.583e-05 & 1.162\\
Balanced 11 & 5.229e-06 & 4.854e-06 & -1.09e-06 & -1.465e-06 & 3.743e-07\\
Near 12 & 9.621 & 7.635 & 1.986 & 6.886e-07 & 1.986\\
Balanced 12 & -4.214e-07 & -2.999e-07 & -7.315e-07 & -6.1e-07 & -1.216e-07\\
Near 13 & 2.773 & 2.773 & 0.0002761 & 1.687e-06 & 0.0002744\\
Balanced 13 & 3.117e-05 & 4.6e-05 & -6.103e-05 & -4.62e-05 & -1.483e-05\\
\bottomrule
\end{tabular}
\end{table}
Reduction in task-distance penalty with nearby header: 7.6266 nats [6.9081, 8.3015].
Reduction in task-distance penalty with distant header: 8.6760 nats [7.7375, 9.6152].
The largest absolute mean co-located-versus-legacy discrepancy is 4.09 nats; these controls quantify the additional distractor regrouping.
\paragraph{Scope and resources.}
This test uses fresh prompts within the existing catalog and existing training lineage. The task-bearing block includes mapping and output-order information, so the contrast does not isolate factual retention from every other task requirement. No new latency certificate, independent training replication or natural-language result is claimed.
The six completed evaluations and one stopped startup attempt use 6.60 additional H100-hours. All are terminal; peak reservation is 48 GPUs within the primary32 plus extra16 allocation. The stopped attempt produced no predictions and was manually replaced with the same frozen evaluation; its scheduler reports zero runtime despite startup logs.
Across completed jobs, whole-job mean GPU busy is 55.1--80.7\%, and memory occupancy is 6.6--8.1\%. GPU busy is not achieved FLOP utilization. These measurements include startup, numerical screens and waiting for slower ranks; they do not satisfy an 80\% utilization target.


\paragraph{What the factorial test identifies.}
Because $D$ is fixed, these layout contrasts are changes in learned conditional
error, not changes in the target's dependence structure. The positive primary
contrast shows that the training-arm advantage survives with nearby
generic instructions; it does not imply that header position has no effect.
Fixed token multisets control content and length, but moving blocks necessarily
changes sequence order. In particular, the co-located versus legacy contrasts
measure sensitivity to regrouping irrelevant records even with the same
meaningful-block positions. We therefore interpret the result as a test of
competence under the declared layouts, not proof of a unique memory or retrieval
circuit. It neither separates all components of the earlier training treatment
nor adds a new architecture-level cost comparison.
